% Self-contained inclusion fragment; requires amsmath, amssymb, amsthm,
% and theorem/proposition environments in the enclosing document.
% New scoped observation. No checker-size or full bi-interpretation claim.
\section{An exact selector boundary for the tied-seed recursor}
\label{sec:selector-repair}

The negative result has a precise positive counterpart. A bounded equality
characteristic cannot remove the numerical obstruction, whereas an unbounded
zero selector gives an immediate compiler for every fixed primitive-recursive
description. The latter compiler is a numerical result. It does not by itself
establish a smaller complete calculus or discharge the full logical,
equality-reflection, induction, and bi-interpretation obligations.

\subsection{Bounded-valued primitives preserve the obstruction}

Consider the grammar with zero and
\[
 B(0;u,v.h)=S\bigl(h(0,0)\bigr),\qquad
 B(Sn;u,v.h)=h\bigl(n,B(n;u,v.h)\bigr),
 \qquad Sa:=B(0;u,v.a),
\]
with the same passive-parameter and capture-free substitution conventions as
before. Extend it by any collection of fixed total operations $e$ such that,
for each operation separately, there is a constant $K_e$ with
\[
                         e(\vec x)\le K_e
                         \quad\hbox{for every input }\vec x.
\]
Only finitely many operations occur in any individual term.

\begin{proposition}[Bounded extensions remain almost unary]
Every term of this enlarged grammar is bounded either by a constant or by
a nondecreasing unary primitive-recursive function of one of its free
coordinates. In particular, adjoining a primitive equality characteristic
with values in $\{0,1\}$ does not invalidate the coordinate-bound obstruction.
\end{proposition}
\begin{proof}
Repeat the structural coordinate-bound argument. A term whose outermost
operation is an added $e$ has the constant bound $K_e$, independently of
the bounds of its arguments. This is the only new structural case. At a
$B$-term the update is again either uniformly bounded or bounded in its
counter coordinate, state coordinate, or one fixed passive coordinate.
The same constant, one-coordinate, and iterated-unary bounds therefore apply.
Thus the finite-dimensional quotient obstruction, whenever its interpreted
operations are required to have term realizers, applies unchanged.
\end{proof}

The proposition is about term operations in the standard algebra. Further
axioms that are sound in that algebra cannot change the bound. It does not
exclude interpretations using formula-defined functional graphs.

\subsection{One unbounded zero selector restores ordinary recursion}

Instead adjoin the numerical operation $C$ with precisely the computation
schemas
\[
                  C(0,a,b)=a,\qquad C(Sn,a,b)=b.
                  \tag{C}
\]
Allow $C$ in subsequent finite $B$ update descriptions. This syntactic
extension is part of the proposed repair and is not implicit. The selector
is unbounded: its two output arguments are arbitrary numerical terms.
It has no one-coordinate bound, since at control zero its output can grow
with $a$, and at positive control its output can grow independently with $b$.

Define
\[
                         W(n):=B(n;u,v.u).
\]
The two $B$ computations give $W(0)=S0$ and $W(Sn)=n$. Let fixed earlier
terms $g(\vec p)$ and $h(n,z,\vec p)$ be given. Choose fresh $u,v$ and put
\[
 \begin{aligned}
 H(u,v,\vec p)&:=C\bigl(u,g(\vec p),h(W(u),v,\vec p)\bigr),\\
 J(n,\vec p)&:=B\bigl(Sn;u,v.H(u,v,\vec p)\bigr).
 \end{aligned}
 \tag{Selector compiler}
\]

\begin{proposition}[Direct primitive-recursion compiler]
Ordinary equality reasoning and the displayed computation schemas prove
\[
 J(0,\vec p)=g(\vec p),\qquad
 J(Sn,\vec p)=h\bigl(n,J(n,\vec p),\vec p\bigr).
\]
No equation induction, predecessor search, equality characteristic,
pairing construction, or independent initializer is used in this derivation.
\end{proposition}
\begin{proof}
Suppress the fixed passive parameters. The first actual update discards
the tied initial state:
\[
 \begin{aligned}
 J(0)
 &=B(S0;u,v.H)\\
 &=H\bigl(0,B(0;u,v.H)\bigr)\\
 &=C\bigl(0,g,h(W0,B(0;u,v.H))\bigr)=g.
 \end{aligned}
\]
At a subsequent update the counter is a displayed successor, so
\[
 \begin{aligned}
 J(Sn)
 &=B(S(Sn);u,v.H)\\
 &=H\bigl(Sn,B(Sn;u,v.H)\bigr)\\
 &=C\bigl(Sn,g,h(W(Sn),J(n))\bigr)\\
 &=h\bigl(n,J(n)\bigr).
 \end{aligned}
\]
These are identities with arbitrary numerical parameters, not a computation
restricted to standard numerals. The original zero value is $B(0;u,v.H)=Sg$;
the compiler deliberately skips it and starts at $B(S0;u,v.H)=g$.
\end{proof}

Zero, successor, and projections are already terms, and composition is
substitution. External structural recursion on a fixed primitive-recursive
description therefore compiles every primitive-recursive function into this
enlarged grammar. Conversely each fixed $B,C$ description has its usual
primitive-recursive numerical realization. These two facts establish
numerical primitive-recursive completeness. The general logical and
interpretation-composite claims require their own proofs and size accounting.

\subsection{Why a formula definition alone does not supply this compiler}

In ordinary intuitionistic logic, suppose numerical coverage and $0\ne S0$
are available. Put $Z(n)=B(n;u,v.0)$. Then $Z0=S0$ and $Z(Sn)=0$;
coverage and these equations prove $z=0\vee z\ne0$. Consequently
\[
 R_C(z,a,b,w):=(z=0\wedge w=a)\vee(z\ne0\wedge w=b)
\]
is provably total and single-valued. Adjoining a function symbol with this
graph is an ordinary definitional extension of that first-order theory.

Two further changes would nevertheless need justification. First, the old
$B$ signature has symbols only for its original term updates. Adjoining
$C$ by a formula definition does not automatically provide a $B$ symbol
for an update containing $C$, or prove totality and uniqueness of that new
recursion graph. Second, literal-equation induction for new $C$-terms
translates into induction on formulas that need not be literal equations
of the old language. Thus closing both the recursor-description scheme and
the induction scheme under such definitions is additional mathematics,
not a free consequence of conservativity of a single function definition.
The explicit positive repair above includes $C$ in the numerical grammar;
it makes no claim that the original tied-seed theory already contains it.
